\section{Evaluation}\label{sec:evaluation}
Ziel der Evaluation ist es, zu untersuchen, ob der implementierte Prototyp die zentralen Prinzipien von ASBAC für Dateiöffnungen unter Linux wie vorgesehen umsetzt. Dazu werden die Verarbeitung von Policies und dynamischen Attributen, die Kontrolle neuer Dateiöffnungen sowie die Nachbewertung bereits bestehender Dateizugriffe überprüft. Zusätzlich werden der Laufzeitaufwand und die technischen Grenzen der prototypischen Umsetzung betrachtet.


\subsection{Testumgebung}

Die Evaluation erfolgt in einer ausschließlich für die Untersuchung verwendeten virtuellen
Maschine. Dadurch lassen sich die für das Laden der eBPF-Programme erforderlichen privilegierten
Operationen ausführen, ohne ein Produktivsystem zu beeinflussen. Während der privilegierten
End-to-End-Tests läuft keine weitere Instanz des Prototyps, da die Tests vorübergehend die
gepinnten BPF-Maps unter \mbox{\path{/sys/fs/bpf/tails-pdp}} verwenden. Die VM-Konfiguration, die
NixOS-Systemkonfiguration und die Definition der Entwicklungsumgebung werden gemeinsam mit dem
Quellcode archiviert. Die Dateien \path{nix-config/configuration.nix},
\path{nix-config/hardware-configuration.nix} und \path{nix-config/shell.nix} dokumentieren die
für die Evaluation relevanten System-, Hardwareerkennungs- und Werkzeugvoraussetzungen.

Für die Kernel- und Durchsetzungstests muss BPF in der aktiven LSM-Reihenfolge enthalten sein.
Darüber hinaus werden BTF-Informationen unter \mbox{\path{/sys/kernel/btf/vmlinux}} und ein unter
\mbox{\path{/sys/fs/bpf}}  eingehängtes BPF-Dateisystem vorausgesetzt. Das Laden und Anheften der
eBPF-Programme, der Zugriff auf die gepinnten Maps und der \texttt{ptrace}-basierte Entzug eines
File Descriptors benötigen erhöhte Berechtigungen. Root-Berechtigung und BTF-Verfügbarkeit werden
vom Testskript explizit geprüft; die Verfügbarkeit von BPF-LSM und BPF-Dateisystem wird durch den
erfolgreichen Runtime-Start, das Anhängen des Hooks und die erzeugten Map-Pins nachgewiesen.

\subsubsection{Hard- und Softwareumgebung}

Die virtuelle Maschine wird mittels QEMU/KVM auf einem Proxmox-VE-Host betrieben. Für die
Reproduzierbarkeit sind dabei zwei Perspektiven zu unterscheiden: Die von Proxmox ausgegebene
VM-Konfiguration beschreibt die vom Hypervisor bereitgestellten Ressourcen, während die von
NixOS erzeugte \path{hardware-configuration.nix} die im Gastsystem erkannten Geräte und die zu
ladenden Kernelmodule festhält. Keine der beiden Konfigurationen ersetzt daher vollständig die
jeweils andere.

Als virtuelles Prozessormodell ist \texttt{x86-64-v2-AES} konfiguriert. Zwei virtuelle Sockel mit
jeweils zwei Kernen stellen insgesamt vier virtuelle CPUs bereit; zusätzlich sind 8192~MiB
Arbeitsspeicher zugewiesen. Das virtuelle Systemlaufwerk besitzt eine Kapazität von 64~GiB und ist
über VirtIO-SCSI angebunden. Diese Angaben werden aus der VM-Konfiguration übernommen, weil sie
nicht Bestandteil der NixOS-Hardwarekonfiguration sind. Weitere betriebliche Kennungen und
Gerätedetails der VM sind für die durchgeführten Funktions- und Laufzeittests nicht
entscheidungsrelevant und werden daher nicht einzeln aufgeführt.

Als Gastbetriebssystem kommt NixOS zum Einsatz. Die archivierte Systemkonfiguration setzt
\texttt{system.stateVersion} auf \texttt{25.05}. Diese Option steuert kompatibilitätsrelevante
Voreinstellungen, um bestehende Anwendungsdaten bei Systemaktualisierungen weiterverwenden zu
können. Sie bezeichnet daher nicht die tatsächlich installierte NixOS-Version
\citep{nixosStateVersion}. Der Kernel wird durch die Konfiguration aus den Quellen des stabilen
Linux-Kernels 6.16.12 mit festgelegtem SHA-256-Hash gebaut. Auf dem Testsystem ist NixOS
\texttt{25.05.810656.5b5be50345d4} installiert. Die Konfiguration aktiviert unter
anderem die Kerneloptionen für das Security-Framework und SecurityFS. Die LSM-Auswahl wird
zentral über die NixOS-Option \texttt{security.lsm} auf \texttt{landlock}, \texttt{yama} und
\texttt{bpf} festgelegt. NixOS erzeugt daraus den Kernelparameter
\texttt{lsm=landlock,yama,bpf}. Zusammen mit dem vom Kernel ergänzten Capability-LSM ergibt
sich die auf dem Testsystem über \path{/sys/kernel/security/lsm} bestätigte aktive Reihenfolge
\texttt{capability,landlock,yama,bpf}.

Die von NixOS generierte \path{hardware-configuration.nix} bindet das QEMU-Gastprofil ein und
legt \texttt{x86\_64-linux} als Plattform des Gastsystems fest. Für den Systemstart werden unter
anderem die Module für VirtIO-PCI, VirtIO-SCSI und SCSI-Datenträger in die Initial Ramdisk
aufgenommen. Das Root-Dateisystem verwendet ext4, die EFI-Systempartition vfat.

Die Entwicklungs- und Build-Umgebung wird mit \texttt{nix-shell} aus
\path{nix-config/shell.nix} erzeugt. Für den regulären Projektbuild sind insbesondere Rustup und
Cargo, LLVM~21 einschließlich seiner Laufzeitbibliothek sowie \texttt{bpf-linker} erforderlich.
Der Shell-Hook wählt die Rust-Nightly-Toolchain \texttt{nightly-2025-08-01} einschließlich
\texttt{rust-src}, \texttt{rustfmt} und Clippy aus. Aya führt den Build der eBPF-Crate für das
Ziel \texttt{bpfel-unknown-none} aus. Das Build-Skript dieser Crate setzt die Verfügbarkeit des
\texttt{bpf-linker} voraus; die Shell-Konfiguration sieht dafür Version 0.9.15 mit
LLVM-21-Unterstützung vor. Die Rust-Abhängigkeiten werden durch \path{Cargo.lock} festgelegt; der
verwendete Stand enthält Aya 0.13.2 und \texttt{aya-ebpf} 0.1.2.

Für die Ausführung der Skripte enthält die Umgebung außerdem Bash, Coreutils, GNU Grep, Procps und
Python~3.

Tabelle~\ref{tab:testumgebung} fasst die aus den archivierten Konfigurationsdateien ableitbaren
Merkmale sowie die auf dem Testsystem ermittelte NixOS-Version und aktive LSM-Reihenfolge
zusammen. Die konkrete Proxmox-, NixOS-, Nixpkgs-, Rust- und Cargo-Version sowie der
getestete Git-Commit werden unmittelbar auf dem Testsystem erfasst und zusammen mit den
Testergebnissen dokumentiert. Dies ist insbesondere für Nixpkgs erforderlich, weil die aktuelle
\path{nix-config/shell.nix} das zum Ausführungszeitpunkt konfigurierte \texttt{<nixpkgs>}
importiert und keine feste Nixpkgs-Revision enthält.

\begin{table}[htbp]
    \centering
    \caption{Hard- und Softwarekonfiguration der Testumgebung}
    \label{tab:testumgebung}
    \begin{tabular}{p{0.36\textwidth}p{0.55\textwidth}}
        \hline
        \textbf{Komponente} & \textbf{Konfiguration} \\
        \hline
        Virtualisierungsplattform & Proxmox VE (9.2.18) mit QEMU/KVM \\
        Virtuelles CPU-Modell & \texttt{x86-64-v2-AES} \\
        CPU-Konfiguration & 4 vCPUs \\
        Arbeitsspeicher & 8192 MiB \\
        Virtuelles Laufwerk & 64 GiB \\
        Speicheranbindung & VirtIO-SCSI \\
        NixOS-Hardwareprofil & QEMU-Gast, Plattform \texttt{x86\_64-linux} \\
        Dateisysteme & ext4 für \path{/}, vfat für \path{/boot} \\
        Gastbetriebssystem & NixOS \texttt{25.05.810656.5b5be50345d4} \\
        Linux-Kernel & 6.16.12, aus festgelegtem Quellarchiv gebaut \\
        LSM-Reihenfolge & \texttt{capability,landlock,yama,bpf} \\
        Entwicklungsumgebung & \texttt{nix-shell} auf Basis von \texttt{<nixpkgs>} \\
        LLVM & LLVM-Paketfamilie 21 einschließlich \texttt{libLLVM} \\
        Rust-Toolchain & \texttt{nightly-2025-08-01} mit \texttt{rust-src} \\
        eBPF-Zielarchitektur & \texttt{bpfel-unknown-none} \\
        \texttt{bpf-linker} & 0.9.15 \\
        Aya / \texttt{aya-ebpf} & 0.13.2 / 0.1.2 gemäß \path{Cargo.lock} \\
        \hline
    \end{tabular}
\end{table}

Da die Evaluation innerhalb einer virtuellen Maschine erfolgt, können insbesondere Zeitmessungen
durch das Scheduling des Hypervisors und durch gleichzeitig auf dem Proxmox-Host ausgeführte
Arbeitslasten beeinflusst werden. Gemessene Laufzeiten charakterisieren daher die dokumentierte
Testumgebung und sind nicht ohne Weiteres als allgemeingültige Werte für andere Hard- und
Softwarekonfigurationen zu verstehen.

\subsection{Teststrategie und Testszenarien}

Die Evaluation kombiniert automatisierte Unit- und Komponententests mit privilegierten
End-to-End-Tests und ausgewählten Laufzeitmessungen. Die Teststufen liefern unterschiedliche
Arten von Evidenz und werden deshalb getrennt betrachtet. Unit- und Komponententests prüfen
isolierbare Logik unter kontrollierten Bedingungen. End-to-End-Szenarien untersuchen dagegen das
Zusammenspiel der Userspace-Komponenten mit den BPF-Maps und dem tatsächlich geladenen
eBPF-LSM-Programm. Die Laufzeitmessungen ergänzen die funktionalen Nachweise um quantitative
Aussagen zur Reaktionszeit und zum zusätzlichen Aufwand bei kontrollierten Dateiöffnungen.

Dieser Abschnitt beschreibt zunächst den Versuchsaufbau und das jeweils erwartete Verhalten. Die
tatsächlich beobachteten Resultate werden davon getrennt in Abschnitt~6.3 dargestellt. Dadurch
wird vermieden, die Bewertungskriterien nach Kenntnis der Ergebnisse anzupassen.

\subsubsection{Teststufen und Bewertungskriterien}

Die Unit-Tests untersuchen einzelne Funktionen mit fest vorgegebenen Eingaben und vergleichen deren
Ausgabe durch Assertions mit dem erwarteten Ergebnis. Komponententests beziehen mehrere Funktionen
einer Userspace-Komponente ein, ersetzen externe Abhängigkeiten jedoch teilweise durch
Testdoubles. So prüfen die Generationstests die Reihenfolge zwischen dem Schreiben in die inaktive
Bank und der anschließenden Aktivierung mit einem speicherbasierten Ersatz für die BPF-Maps.
Die Tests des Userspace-PEP verwenden einen
\texttt{FakeFdCloser}, der Schließversuche protokolliert, aber keinen realen Systemaufruf in einem
anderen Prozess ausführt. Diese Tests können die fachliche Logik und das lokale Fehlerverhalten
präzise prüfen, belegen jedoch weder die Annahme des eBPF-Programms durch den Kernel-Verifier noch
die reale Durchsetzung einer Zugriffsentscheidung.

Die End-to-End-Tests führen den vollständigen Pfad von einer Policy- oder Attributdatei über die
Loader und gepinnten BPF-Maps bis zur beobachtbaren Wirkung auf einen Dateizugriff aus. Sie prüfen
damit Eigenschaften, die sich in einer isolierten Testumgebung nicht nachweisen lassen, etwa das
Laden des eBPF-Objekts, das Anhängen des \texttt{file\_open}-Hooks und den tatsächlichen Entzug
eines File Descriptors. Quantitative Messungen werden wiederholt ausgeführt und anhand vorab
festgelegter Messgrößen ausgewertet. Einzelne erfolgreiche Durchläufe werden nicht als Nachweis
allgemeiner zeitlicher Garantien oder unbegrenzter Skalierbarkeit interpretiert.

Tabelle~\ref{tab:teststufen} fasst die Teststufen und ihre Aussagegrenzen zusammen.

\begin{table}[htbp]
    \centering
    \caption{Teststufen der Evaluation}
    \label{tab:teststufen}
    \begin{tabular}{p{0.22\textwidth}p{0.32\textwidth}p{0.37\textwidth}}
        \hline
        \textbf{Teststufe} & \textbf{Gegenstand} & \textbf{Aussagegrenze} \\
        \hline
        Unit-Test & Einzelne Parser-, Vergleichs- und Hilfsfunktionen & Kein realer Kernel- oder
        Durchsetzungsnachweis \\
        Komponententest & Zusammenwirkende Funktionen einer Userspace-Komponente, teilweise mit
        Testdoubles & Externe Schnittstellen wie BPF-Maps oder \texttt{ptrace} werden teilweise
        ersetzt \\
        End-to-End-Test & Vollständiger Pfad auf dem Linux-Zielsystem & Nachweis gilt für das
        definierte Szenario und die dokumentierte Testumgebung \\
        Laufzeitmessung & Reaktionszeit und Aufwand kontrollierter Dateiöffnungen & Keine
        allgemeine Echtzeit- oder Skalierbarkeitsgarantie \\
        \hline
    \end{tabular}
\end{table}

Für die abschließende Zuordnung zu den Anforderungen werden die Bewertungen \emph{erfüllt},
\emph{teilweise erfüllt} und \emph{nicht erfüllt} verwendet. Eine
Anforderung gilt als erfüllt, wenn ihr Akzeptanzkriterium unter den dokumentierten Bedingungen
vollständig beobachtet wird. Als teilweise erfüllt wird sie bewertet, wenn das vorgesehene
Grundverhalten nachweisbar ist, aber relevante Einschränkungen oder nicht abgedeckte Fälle
verbleiben. Schlägt das erwartete Verhalten fehl, gilt die Anforderung als nicht erfüllt.
Die Bewertung bezieht sich
stets auf den begrenzten Funktionsumfang des Prototyps und stellt keinen vollständigen
Sicherheitsnachweis dar.

\subsubsection{Unit- und Komponententests}

Das Skript \path{test.sh} bildet den zentralen Einstiegspunkt der gesamten Prüfkette. Sein erster,
nicht privilegierter Teil führt 61 automatisierte Rust-Tests in fünf Testbereichen aus. Die Option
\texttt{--locked} stellt sicher, dass Cargo dabei die
in \path{Cargo.lock} festgehaltene Abhängigkeitsauflösung nicht verändert. Der in der
Cargo-Konfiguration für reguläre Programmstarts eingestellte privilegierte Runner wird für diese
Tests durch \texttt{env} ersetzt. Die Tests benötigen daher keine Root-Rechte, laden kein
eBPF-Programm in den Kernel und verändern keine gepinnten BPF-Maps.

Die Tests der gemeinsamen Crate \path{tails-pdp-common} bilden mit 16 Testfällen die fachliche
Entscheidungslogik ab. Sie prüfen das Matching statischer und zeitabhängiger Policies, die
Ableitung von UTC-Komponenten, die Grenzen der Vergleichsoperatoren und die Combining-Regel
\emph{deny-overrides}. Darüber hinaus untersuchen sie dynamische Zahlen-, Boolean- und
Stringattribute, Typabweichungen sowie die unterschiedlichen Objektidentifikatoren der System-,
Subject- und Resource-Namespaces. Da Teile dieser Funktionen sowohl im Kernelspace- als auch im
Userspace-Pfad verwendet werden, prüfen die Tests die gemeinsam genutzte Semantik, nicht jedoch
die Einbindung in die jeweilige Laufzeitumgebung.

Der Policyloader besitzt 16 Tests für das Einlesen, Parsen, Übersetzen und Validieren textueller
Policies. Abgedeckt sind unter anderem statische und zeitabhängige Policies, frei benannte
dynamische Attribute, ungültige Attributnamen, doppelte Policynamen, fehlende Semikolons,
unbekannte oder deaktivierte Actions, ungültige Zeitwerte und das Überschreiten der festgelegten
Kapazitätsgrenzen. Zwei Tests prüfen für statische und dynamische Policies die Annahme von
Kommandonamen mit 15 Bytes und die Ablehnung von 16 Bytes, jeweils auch mit Zeichen, deren
UTF-8-Kodierung mehrere Bytes benötigt. Weitere Tests prüfen das rekursive Einlesen ausschließlich von Dateien mit der
Endung \path{.policy}. Ein speicherbasierter \texttt{FakePolicyStore} kontrolliert, dass zuerst
die inaktive Bank beschrieben und erst danach die neue Generation aktiviert wird. Simulierte
Schreib- und Aktivierungsfehler müssen die zuvor aktive Generation unverändert lassen; außerdem
werden nicht mehr belegte Einträge einer neu vorbereiteten Bank deaktiviert.

Die 13 Tests des Attributloaders prüfen das Parsen der unterstützten Zahlen-, Boolean- und
Stringwerte sowie die Validierung von Attributnamen und des zulässigen DEFCON-Wertebereichs.
Ergänzend werden ein vollständiges Verzeichnis mit System-, Subject- und Resource-Attributen sowie
die transaktionale Aktivierung einer neuen Attributbank mit simulierten Lösch-, Schreib- und
Aktivierungsfehlern untersucht. Weitere Fehlerfälle decken die Kapazitätsprüfung vor der ersten
Map-Änderung, Fehler beim Lesen von Generation und Belegung sowie die Fortsetzung des laufenden
Updaters nach einer fehlgeschlagenen Übernahme eines neuen Attributstands ab. Zwei
weitere Tests in \path{tails-pdp-userspace-common} untersuchen den begrenzten Trigger-Kanal. Sie
prüfen, dass mehrere schnell aufeinanderfolgende Aktivierungen bei voller Kanalkapazität
koalesziert werden und ein geschlossener Kanal als Fehler gemeldet wird.

Die 14 Tests des Userspace-PEP prüfen die Auswertung von Prozess- und
File-Descriptor-Informationen sowie die Steuerung der Nachbewertung. Dazu gehören das Filtern
numerischer Einträge unter \path{/proc}, die Auswahl der Real UID aus
\path{/proc/<pid>/status}, die Erkennung regulärer Dateien und die gezielte Zuordnung eines
Verstoßes zu Prozess und File Descriptor. Mit dem \texttt{FakeFdCloser} wird geprüft, dass ein
File Descriptor innerhalb eines Scans höchstens einmal behandelt wird, unterschiedliche
Deskriptoren unabhängig bleiben und ein fehlgeschlagener Schließversuch die Bearbeitung weiterer
Verstöße nicht abbricht. Ein zusätzlicher Test bildet die Konjunktion mehrerer dynamischer
Bedingungen im produktiven Lookup-Ablauf ab; ein weiterer weist nach, dass ein fehlgeschlagener
Entzug die Verarbeitung eines anderen Zielprozesses nicht verhindert. Ergänzend testen asynchrone
Fälle das Warten auf Aktivierungsereignisse sowie die Berechnung relevanter Stunden- und
Modulo-Zeitgrenzen.

Neben den funktionalen Rust-Tests führt der nicht privilegierte Teil von \path{test.sh} drei
Qualitätsprüfungen aus. Eine reine
Formatprüfung kontrolliert mit \texttt{cargo fmt}, ob der Rust-Quellcode dem vorgegebenen Format
entspricht. Clippy analysiert Produktions- und Testcode statisch; Warnungen werden dabei als
Fehler behandelt. Abschließend erzeugt ein Release-Build die Userspace-Binaries und über das
Build-Skript des Hauptprogramms auch das eingebettete eBPF-Objekt. Dieser Build belegt die
Übersetzbarkeit des gemeinsamen Quellstands, ersetzt aber nicht die anschließende Prüfung durch
den Kernel-Verifier. In der eBPF-Crate selbst werden keine nativen Unit-Tests ausgeführt, da ihr
Code für die BPF-Zielarchitektur und eine \texttt{no\_std}-Umgebung übersetzt wird. Die
gemeinsame Policylogik wird daher nativ getestet; die tatsächliche Kernelausführung ist
Gegenstand der folgenden End-to-End-Szenarien.

Ergänzend prüfen sechs nicht privilegierte Python-Tests die Laufzeit (PERF-03) Messauswertung. Sie prüfen
unter anderem die zeitliche Reihenfolge, fehlende oder mehrfache Zeitmarkierungen, eine
unvollständige Bündelungsphase und noch laufende Scans vor Messbeginn.

\subsubsection{Privilegierte End-to-End-Tests}

Die End-to-End-Prüfung wird auf dem in Abschnitt~6.1 beschriebenen Linux-Zielsystem mit
Root-Rechten ausgeführt. Der zentrale Runner \path{test.sh} fordert diese Berechtigungen nach dem
Release-Build über \texttt{sudo} an und ruft \path{tests/test-e2e.sh} auf. Dieses Skript bildet den
gemeinsamen Einstiegspunkt der End-to-End-Stufe und
ruft für die Test-IDs \texttt{E2E-01} bis \texttt{E2E-19} jeweils eine eigene Bashdatei unter
\path{tests/e2e/} auf. Vor dem Start bricht die Hauptdatei ab, wenn bereits eine Instanz des
Prototyps läuft; anschließend entfernt sie die bekannten Test-Maps unter
\path{/sys/fs/bpf/tails-pdp}. Deshalb darf die Suite ausschließlich auf dem dedizierten Testsystem
und nicht parallel zu einer regulären Instanz oder einem weiteren Testlauf ausgeführt werden.

Für die ersten zehn Szenarien startet die Hauptdatei eine gemeinsame Runtime in einem temporären
Arbeitsverzeichnis. Dieses enthält eigene Policy- und Attributverzeichnisse sowie eine geschützte
und eine nicht geschützte Testdatei. Änderungen werden zunächst vollständig in eine temporäre
Datei geschrieben und anschließend atomar aktiviert. Wo ein erlaubter Zugriff allein noch keinen
Nachweis für die Verarbeitung einer Änderung liefern würde, wartet der Test zusätzlich auf einen
Wechsel der Policy- oder Attributgeneration. Nach \texttt{E2E-10} wird diese Runtime beendet und
der gepinnte Testzustand entfernt.

Die Szenarien \texttt{E2E-11} bis \texttt{E2E-19} werden darin anschließend über
\path{tests/test-evaluation.sh} mit jeweils eigener Runtime und isoliertem Arbeitsverzeichnis
ausgeführt. Der zugehörige Runner sperrt parallele Evaluationsläufe, verwaltet ausschließlich die
von ihm gestarteten Prozesse und verweigert den Start, wenn bekannte gepinnte Maps bereits
vorhanden sind. Beim Aufräumen entfernt er nur die durch den jeweiligen Lauf erzeugten Maps. Für
die Prüfung neuer Dateiöffnungen verwendet dieser Teil der Suite einen kleinen Syscall-Helfer. Ein
Zugriff gilt nur dann als durch den LSM-Hook verweigert, wenn \texttt{open()} mit \texttt{EPERM}
fehlschlägt; andere Ein-/Ausgabefehler werden nicht als erfolgreiches Deny gewertet.

Die Kennungen der ergänzenden Szenarien unterscheiden Komponententests (COMP),
Charakterisierungstests (CHAR), Tests auf Wettlaufsituationen (RACE), Laufzeitmessungen (PERF),
Kapazitätstests (LOAD) und Stabilitätstests (STAB). Die nachgestellte Nummer identifiziert das
jeweilige Szenario innerhalb der Testgruppe.

Nach Abschluss aller 19 End-to-End-Szenarien ruft \path{test.sh} denselben Evaluationsrunner ein
zweites Mal ausschließlich für \texttt{COMP-04}, \texttt{CHAR-01}, \texttt{RACE-01},
\texttt{PERF-01} bis \texttt{PERF-03}, \texttt{LOAD-01} und \texttt{STAB-01} auf. Damit werden
alle Teststufen durch einen Aufruf von \path{test.sh} ausgeführt, ohne die Szenarien
\texttt{E2E-11} bis \texttt{E2E-19} doppelt zu starten. Die Teilrunner bleiben für gezielte
Wiederholungen einzelner Szenarien separat aufrufbar.

Tabelle~\ref{tab:e2e-basis} fasst die grundlegenden Start-, Entscheidungs- und Attributszenarien
zusammen. Die Kriterien beschreiben den vorab erwarteten Zustand; die tatsächlich beobachteten
Ergebnisse werden getrennt in Abschnitt~6.3 berichtet.

\begin{table}[htbp]
    \centering
    \caption{Grundlegende End-to-End-Szenarien}
    \label{tab:e2e-basis}
    \begin{tabular}{p{0.16\textwidth}p{0.34\textwidth}p{0.41\textwidth}}
        \hline
        \textbf{Test-ID} & \textbf{Gegenstand} & \textbf{Bestehenskriterium} \\
        \hline
        E2E-01 & Runtime, Kernel-Verifier, LSM-Attach und Maps & Die Runtime bleibt aktiv und alle
        erwarteten Maps sind gepinnt. \\
        E2E-02 bis E2E-04 & Standardentscheidung, statische Deny-Policy und aktuelle UTC-Stunde &
        Der Zugriff ist ohne passende Deny-Policy erlaubt, wird bei passender statischer oder
        zeitabhängiger Policy verweigert und nach deren Entfernung wieder erlaubt. \\
        E2E-05 bis E2E-07 & System-, Subject- und Resource-Attribute & Änderungen von
        \texttt{system.defcon}, \texttt{subject.position} und
        \texttt{resource.classification} verändern die reale Zugriffsentscheidung in beide
        Richtungen. \\
        E2E-08 & Ungültige Policygeneration & Die Runtime meldet den Parserfehler, aktiviert die
        fehlerhafte Generation nicht und erzwingt weiterhin die letzte gültige Deny-Policy. \\
        E2E-12 & Konjunktion mehrerer Attribut-Namespaces & Die Deny-Policy greift nur bei
        gleichzeitig passenden System-, Subject- und Resource-Attributen; fehlende, typfalsche
        oder abweichende Einzelwerte heben das Deny auf. \\
        E2E-13 & Ungültige Attributgeneration & Ein unzulässiger Attributwert wird gemeldet, ohne
        die sichtbare Generation oder die zuvor gültige Entscheidung zu verändern. \\
        \hline
    \end{tabular}
\end{table}

Die weiteren Szenarien in Tabelle~\ref{tab:e2e-enforcement} konzentrieren sich auf die
Combining-Regel, bestehende File Descriptors, die Administrationsschnittstelle und kontrollierte
Fehlerfälle. Beim nachträglichen Entzug hält ein Hilfsprozess mehrere Deskriptoren geöffnet und
prüft über \texttt{fstat()}, ob ausschließlich die von der neuen Entscheidung betroffenen
Deskriptoren ungültig werden. Der Fehlercode \texttt{EBADF} (\emph{Bad file descriptor},
ungültiger Dateideskriptor) zeigt dabei an, dass der geprüfte Deskriptor nicht mehr gültig ist
\citep{manErrno}.

\begin{table}[htbp]
    \centering
    \caption{End-to-End-Szenarien für Enforcement und Fehlerverhalten}
    \label{tab:e2e-enforcement}
    \begin{tabular}{p{0.16\textwidth}p{0.34\textwidth}p{0.41\textwidth}}
        \hline
        \textbf{Test-ID} & \textbf{Gegenstand} & \textbf{Bestehenskriterium} \\
        \hline
        E2E-09 und E2E-16 & Selektiver Entzug bestehender File Descriptors & Ein beziehungsweise
        mehrere verletzende Deskriptoren werden geschlossen; weiterhin zulässige Deskriptoren
        desselben Prozesses bleiben geöffnet. \\
        E2E-10 und E2E-14 & Administrationsschnittstelle und Nachvollziehbarkeit & Die Ausgabe
        enthält aktive Generation, Deny-Policy, Ressourcenidentität, Attributbedingung und
        Attributwert, ohne den Zustand zu verändern. Dies weist die rekonstruierbare Ursache im
        kontrollierten Szenario nach, jedoch kein individuelles Auditlog jeder Entscheidung. \\
        E2E-11 & Combining-Regel \emph{deny-overrides} & Bei gleichzeitig passender Permit- und
        Deny-Policy wird der Zugriff unabhängig von der lexikalischen Policyreihenfolge verweigert;
        nach Entfernen des Deny wird er erlaubt. \\
        E2E-15 & Real und Effective UID & Bei privilegiertem Effective- und Filesystem-User wird
        die Entscheidung anhand der gesetzten Real UID getroffen; eine andere Real UID bildet die
        Negativkontrolle. \\
        E2E-17 & Nicht erreichbarer Zielprozess beim FD-Entzug & Der gezielt erzeugte
        \texttt{ptrace}-Attach-Konflikt wird protokolliert, ein anderer Zielprozess wird dennoch
        verarbeitet und weitere Policyaktivierungen sowie Kernelentscheidungen bleiben möglich. \\
        E2E-18 & Attributgetriggerter FD-Entzug & Eine reine Attributänderung schließt drei
        unzulässige FDs; zwei erlaubte bleiben offen. Die Policygeneration bleibt unverändert. \\
        E2E-19 & Zeitgetriggerter FD-Entzug & Eine Zeitgrenze schließt drei unzulässige FDs;
        zwei erlaubte bleiben offen. Policy- und Attributgeneration bleiben unverändert. \\
        \hline
    \end{tabular}
\end{table}

Ein Szenario gilt nur dann als bestanden, wenn sein Skript ohne fehlgeschlagene Assertion endet,
die Runtime während der vorgesehenen Prüfung aktiv bleibt und im neu hinzugekommenen Kernel-Log
kein Hinweis auf \emph{Kernel panic}, \emph{Oops}, \emph{BUG} oder einen General-Protection-Fault
enthalten ist. Für die isolierten Szenarien werden Status, Parameter und gegebenenfalls erzeugte
Detaildaten in \path{report.json} gespeichert. Runtime-, Szenario- und Kernel-Logs bleiben im
ausgegebenen Verzeichnis unter \path{/tmp/tails-eval-*} erhalten. Diese Artefakte bilden die
Grundlage für die Ergebnisdarstellung, sind aber selbst noch kein Nachweis über andere Kernel,
Konfigurationen oder Lastsituationen hinaus.
Die Nachweise der hier ausgewerteten Prüfkette sind einschließlich Gesamtprotokoll,
Szenarioberichten, Messrohwerten und Umgebungsangaben unter
\path{thesis/test-results/final-evaluation/} archiviert.

\clearpage
\subsection{Ergebnisse}

Die nachfolgend dargestellten Ergebnisse wurden auf dem in Abschnitt~6.1 beschriebenen
Zielsystem erhoben.

\subsubsection{Ergebnisse der funktionalen Prüfungen}

Alle 61 Unit- und Komponententests wurden erfolgreich ausgeführt. Davon entfallen 16 Tests auf
die gemeinsame Policylogik, 16 auf den Policyloader, 13 auf den Attributloader, zwei auf den
gemeinsamen Triggerkanal und 14 auf den Userspace-PEP. Die Formatprüfung und der separate
Release-Build der Runtime und des Administrationstools einschließlich des eBPF-Objekts waren
ebenfalls erfolgreich. Damit ist nachgewiesen, dass der Prototyp übersetzbar ist
und die in Abschnitt~6.2.2 beschriebenen Assertions erfüllt.

Die vollständige Ausführung von \path{test.sh} wurde mit Exitstatus null abgeschlossen.
Clippy prüfte Produktions-, Test- und Binärziele mit als Fehler behandelten Warnungen erfolgreich.
Auch die sechs ergänzenden Python-Tests zur Auswertung der PERF-03-Zeitmarkierungen bestanden.
Alle acht ergänzenden Evaluationsszenarien für Map-Zugriff, Deskriptorgrenzen, Race-Verhalten,
Performance, Kapazität und Stabilität wurden ebenfalls erfolgreich abgeschlossen.
Damit waren sowohl die funktionalen Prüfungen als auch die definierte Qualitätsprüfkette
vollständig erfolgreich.

Auf dem Zielkernel bestanden alle 19 End-to-End-Szenarien. Die Runtime wurde vom
Kernel-Verifier akzeptiert, der \texttt{file\_open}-Hook angehängt und die erwarteten Maps wurden
gepinnt. Ohne passende Deny-Policy blieb der Zugriff erlaubt; statische, zeitabhängige und von
dynamischen Attributen abhängige Deny-Policies führten dagegen zu \texttt{EPERM}. Änderungen von
System-, Subject- und Resource-Attributen wurden in beide Entscheidungsrichtungen übernommen.
Auch die Konjunktion mehrerer Attribut-Namespaces und die Combining-Regel
\emph{deny-overrides} zeigten das erwartete Verhalten.

Ungültige Policy- und Attributstände ersetzten die jeweils letzte gültige Generation nicht.
Nach der Korrektur des Attributloaders wurde auch eine Überschreitung der realen Map-Kapazität
kontrolliert abgelehnt: Je 512 Attribute pro vollständig belegter Bank wurden akzeptiert, ein
Update auf 513 Attribute bei weiterhin 512 Einträgen der aktiven Bank dagegen vor einer
Veränderung der Maps zurückgewiesen. Generation und aktive Attributwerte blieben unverändert;
anschließende gültige Updates wurden ohne Neustart der Runtime verarbeitet. Für Policies wurden
je 16 statische und 16 Stream-Policies akzeptiert und der jeweils 17. Eintrag kontrolliert
abgelehnt. Damit wurde nicht nur das Verhalten an den vorgesehenen Grenzen, sondern auch die
Wiederherstellung des regulären Updatebetriebs nach einer Ablehnung beobachtet.

Der nachträgliche Entzug bestehender Dateizugriffe war in den vorgesehenen Szenarien ebenfalls
erfolgreich. Sowohl ein einzelner als auch drei gleichzeitig verletzende File Descriptors wurden
geschlossen, während die erlaubten Deskriptoren desselben Prozesses geöffnet blieben.
E2E-18 bestätigte dieses Verhalten nach einer reinen Attributänderung bei unveränderter
Policygeneration. E2E-19 bestätigte es beim Erreichen einer Zeitgrenze ohne Änderung der
Policy- und Attributgeneration. In beiden Szenarien wurden drei unzulässige Deskriptoren
geschlossen, während zwei erlaubte geöffnet blieben. War ein
Zielprozess aufgrund eines bestehenden \texttt{ptrace}-Attachs nicht erreichbar, wurde der
Fehler protokolliert und ein weiterer Prozess trotzdem erfolgreich verarbeitet. Nachfolgende
Policyaktivierungen und Kernelentscheidungen blieben möglich. Die Administrationsschnittstelle
zeigte aktive Generationen, Policies, Ressourcen und Attribute ohne beobachtete Zustandsänderung.
Im kontrollierten Szenario ließ sich die Ursache einer Ablehnung aus der einzigen aktiven
Deny-Policy, der Ressourcenidentität, der Attributbedingung und dem aktuellen Attributwert
rekonstruieren. Ein individuelles Entscheidungsprotokoll für jede Dateiöffnung oder der
ursprüngliche Policyname wird jedoch nicht ausgegeben.

\begin{table}[htbp]
    \centering
    \caption{Zusammenfassung der funktionalen Testergebnisse}
    \label{tab:funktionale-ergebnisse}
    \begin{tabular}{p{0.28\textwidth}p{0.18\textwidth}p{0.45\textwidth}}
        \hline
        \textbf{Prüfung} & \textbf{Ergebnis} & \textbf{Einordnung} \\
        \hline
        Rust-Tests & 61 von 61 bestanden & Unit- und Komponentennachweise ohne realen
        Kernelpfad \\
        End-to-End-Tests & 19 von 19 bestanden & Vollständiger Pfad einschließlich
        Kernel-Verifier, LSM-Attach und Enforcement \\
        Formatprüfung & Bestanden & Keine Formatabweichungen im Rust-Code \\
        Python-Tests & 6 von 6 bestanden & Prüfung der PERF-03-Messauswertung \\
        Clippy-Gesamtprüfung & Bestanden & Statische Analyse mit als Fehler behandelten
        Warnungen \\
        Release-Build & Bestanden & Userspace-Binaries und eBPF-Objekt wurden erzeugt \\
        Kapazitätstest & Bestanden & Kontrollierte Ablehnung ohne Zustandsverlust und mit
        erfolgreichen Folgeupdates \\
        \hline
    \end{tabular}
\end{table}

\subsubsection{Robustheit und Grenzen des Entzugsmechanismus}

Im Stabilitätstest wurden 100 aufeinanderfolgende Zyklen mit gültigen und ungültigen Policy- und
Attributänderungen ausgeführt. Die Runtime blieb bis zum regulären Testende aktiv. Die vom Runner
ausgewertete Differenz des Kernel-Logs enthielt keinen Hinweis auf Kernel-Panic, Oops, BUG oder
einen General-Protection-Fault. Der Abschlusslauf benötigte 175~Sekunden. Dieses Ergebnis
belegt das Verhalten während des begrenzten Tests, stellt jedoch weder einen Langzeittest noch
einen allgemeinen Nachweis der Kernelstabilität dar.

Die Charakterisierung duplizierter und durch \texttt{fork()} vererbter Deskriptoren zeigte, dass
in Eltern- und Kindprozess jeweils alle vier betroffenen Deskriptoren geschlossen wurden, während
jeweils zwei erlaubte Deskriptoren geöffnet blieben. Eine bereits bestehende
\texttt{mmap()}-Abbildung der Datei blieb danach lesbar. Der Prototyp entzieht somit den File
Descriptor, widerruft aber keinen zuvor eingerichteten speicherabgebildeten Zugriff. Diese Grenze
folgt aus dem in Kapitel~3 ausdrücklich auf File Descriptors eingegrenzten Prototypumfang, ist
für die sicherheitstechnische Interpretation des Ergebnisses jedoch relevant.

Im nebenläufigen Stresstest wurden während zehn Policyzyklen 5387 Wechsel der Ressource hinter
derselben Deskriptornummer erzeugt. Dabei wurde kein erlaubter Deskriptor fälschlich geschlossen.
Da der Test stochastisch ist, beweist dieses Ergebnis keine Race-Freiheit. Zwischen dem Erfassen
der Deskriptoridentität und dem späteren \texttt{ptrace}-basierten Schließen verbleibt grundsätzlich
ein Zeitfenster, in dem der Zielprozess die Deskriptornummer wiederverwenden kann.

\subsubsection{Quantitative Ergebnisse}

Tabelle~\ref{tab:performance-ergebnisse} enthält die zentralen Messwerte. Bei PERF-01 wurden je
20\,000 Dateiöffnungen mit warmem Cache gemessen. Jeder der drei Messzustände wurde zuvor
mit 1000 ungemessenen Öffnungs- und Schließvorgängen vorbereitet. Der Median stieg von
3,512~$\mu$s ohne aktiven Prototyp auf 3,985~$\mu$s mit aktivem Prototyp ohne Policy und auf
4,014~$\mu$s bei passender Permit-Policy. Dies entspricht einem medianen Mehraufwand von
13,5 beziehungsweise 14,3~Prozent. Die Messung umfasst \texttt{os.open()} sowie den Aufwand der
Python-Laufzeit und des Timers; \texttt{close()} liegt außerhalb des Messintervalls. Sie isoliert
daher nicht ausschließlich die Ausführungszeit des LSM-Hooks.

\begin{table}[htbp]
    \centering
    \caption{Quantitative Messergebnisse auf dem Zielsystem}
    \label{tab:performance-ergebnisse}
    \begin{tabular}{p{0.40\textwidth}p{0.13\textwidth}p{0.17\textwidth}p{0.17\textwidth}}
        \hline
        \textbf{Messgröße} & \textbf{$n$} & \textbf{Median} & \textbf{p95} \\
        \hline
        Dateiöffnung ohne Runtime & 20\,000 & 3,512~$\mu$s & 3,670~$\mu$s \\
        Dateiöffnung, Runtime ohne Policy & 20\,000 & 3,985~$\mu$s & 4,104~$\mu$s \\
        Dateiöffnung, passende Permit-Policy & 20\,000 & 4,014~$\mu$s & 4,195~$\mu$s \\
        Policyänderung bis neues Deny & 10 & 104,42~ms & 114,75~ms \\
        Attributänderung bis neues Deny & 10 & 99,17~ms & 104,57~ms \\
        Policyänderung bis \texttt{EBADF} & 10 & 110,43~ms & 112,24~ms \\
        \hline
    \end{tabular}
\end{table}

Die Reaktionszeit vom atomaren Austausch einer Datei bis zur beobachtbaren neuen Entscheidung
lag im Median bei 104,42~ms für Policyänderungen und 99,17~ms für Attributänderungen. Dabei wurde
der Zugriff in Intervallen von 10~ms geprüft. Die Messung enthält somit neben Verarbeitung und
Aktivierung auch Dateisystembeobachtung, Scheduling und Polling.

Für den Entzug bestehender Zugriffe betrug die Zeit von der Policydateiänderung bis zum
beobachteten \texttt{EBADF} im Median 110,43~ms. PERF-03 bereitet die Policydatei außerhalb
des überwachten Verzeichnisses vor und beginnt die Messung unmittelbar vor ihrer atomaren
Umbenennung in das Policyverzeichnis. Zuvor müssen der wartende Policy-Watcher und abgeschlossene
Scans über 300~ms einen unveränderten Ruhezustand zeigen. Zeitmarkierungen prüfen, dass das
Dateiereignis erst nach Messbeginn empfangen und die vollständige 100-ms-Bündelungsphase
abgewartet wurde. Zwei weitere Zeitmarkierungen unmittelbar vor und nach dem Map-Systemaufruf
zur Aktivierung grenzen deren Zeitpunkt ein. Die Messung enthält den Aufwand dieser
Instrumentierung; der Entzug wird durch FD-Prüfungen im Abstand von 1~ms beobachtet.
Über zehn Wiederholungen lagen die unteren Grenzen der Aktivierung-bis-Entzug-Intervalle
zwischen 8,09 und 10,22~ms, die oberen Grenzen zwischen 8,10 und 10,25~ms. Eine scheinexakte
Einzelzeit wird daraus nicht abgeleitet. Bei nur zehn Wiederholungen entspricht das angegebene p95 jeweils dem beobachteten
Maximum und ist keine belastbare Schätzung seltener Ausreißer. Da für OA-03 keine feste
Latenzgrenze vorgegeben wurde, bedeutet ein bestandener Performance-Test ausschließlich, dass die
Messung erfolgreich durchgeführt und der Zustandswechsel beobachtet wurde.

\subsubsection{Bewertung der Anforderungen}

Die funktionalen Anforderungen FA-01 bis FA-10 werden innerhalb des in Kapitel~3 definierten
Prototypumfangs als erfüllt bewertet. Die Tests weisen das Einlesen, Aktualisieren und Validieren
von Policies, die Kontrolle neuer Dateiöffnungen, dynamische und frei benannte Attribute sowie
den selektiven Entzug unzulässig gewordener File Descriptors nach. Diese Bewertung umfasst nur
die getesteten Policyformen, Attributtypen und den Hook \texttt{file\_open}.
Bei FA-06 bezieht sich die Erfüllung auf die geforderte Nachbewertung und den gezielten
Schließversuch einschließlich kontrollierter Fehlerbehandlung. E2E-09 und E2E-16 zeigen den
erfolgreichen selektiven Entzug nach einer Policyänderung. E2E-18 und E2E-19 weisen ihn außerdem
nach einer reinen Attributänderung beziehungsweise beim Erreichen einer Zeitgrenze nach.
E2E-17 zeigt, dass ein fehlgeschlagener
Versuch bei einem nicht erreichbaren Zielprozess protokolliert wird und die Behandlung eines
anderen Zielprozesses nicht verhindert. Der im Fehlerfall verbleibende offene Deskriptor ist
damit eine Grenze des Best-Effort-Entzugs, bei dem das Schließen versucht, aber nicht garantiert
wird. Ein im Fehlerfall offen gebliebener Deskriptor gilt nicht als erfolgreich entzogen. Die Bewertung besagt nicht,
dass jeder unzulässig gewordene Zugriff erfolgreich beendet werden kann. Insbesondere folgt
aus FA-06 kein Widerruf bestehender \texttt{mmap()}-Abbildungen, da der bestehende Zugriff dort für
den Prototyp explizit als geöffneter File Descriptor definiert ist. Die beobachtete
Entzugslatenz ist außerdem keine Echtzeitgarantie.

OA-01 wird für die untersuchten Eingaben und die Dauer der Testläufe erfüllt. Ungültige Eingaben,
Kapazitätsüberschreitungen und ein nicht erreichbarer Zielprozess führten in den Abschlussläufen
weder zum Abbruch der Runtime noch zu einem vom Runner erkannten Kernelproblem. Eine allgemeine
Stabilitätsgarantie kann daraus nicht abgeleitet werden. OA-02 ist teilweise erfüllt: Die
Administrationsschnittstelle ermöglicht im kontrollierten Szenario eine Rekonstruktion der
Entscheidungsgrundlage, stellt aber kein individuelles Auditlog bereit und ordnet eine konkrete
Dateiöffnung nicht unmittelbar einem ursprünglichen Policynamen zu. OA-03 ist erfüllt, weil der
zusätzliche Aufwand und die Reaktions- und Entzugslatenz quantitativ untersucht wurden.
OA-04 wird durch die archivierten Konfigurationen,
Szenarioskripte, Testparameter und Rohdaten erfüllt.

Die entwicklungsbezogenen Anforderungen sind nicht allein durch Laufzeittests entscheidbar. Die
Trennung in eBPF-Programm, gemeinsame Datentypen und Entscheidungslogik, Loader,
Userspace-PEP und Administrationstool unterstützt EA-01. Frei benannte Attribute und getrennte
Komponenten bieten Erweiterungspunkte im Sinne von EA-02; neue Attributtypen, Map-Strukturen oder
LSM-Hooks erfordern jedoch weiterhin Änderungen an mehreren Schnittstellen, weshalb diese
Anforderung nur teilweise erfüllt ist. EA-03 wird dadurch unterstützt, dass Einlesen,
Validierung, Aktualisierung, Nachbewertung und Administration im Userspace liegen. Im Kernel
verbleiben der notwendige LSM-Hook, Map-Lookups und die unmittelbar für eine Dateiöffnung
erforderliche Auswertung.

Insgesamt zeigen die Ergebnisse, dass der Prototyp die für diese Arbeit ausgewählten
ASBAC-Prinzipien auf dem untersuchten Linux-System funktional umsetzt: Zugriffsentscheidungen
reagieren auf Policy- und Attributströme, neue Dateiöffnungen werden im Kernel kontrolliert und
bereits geöffnete File Descriptors im Userspace nachbewertet. Die Ergebnisse sind ein
Machbarkeitsnachweis für den definierten Prototyp, kein vollständiger Sicherheits-,
Race-Freiheits-, Echtzeit- oder Produktionstauglichkeitsnachweis.

\newpage
